% !TEX program = xelatex
\documentclass{article}
\usepackage{geometry}\geometry{margin=23mm}
\usepackage{fontspec}\setmainfont[AutoFakeBold=2,AutoFakeSlant=.18]{Songti TC}
\XeTeXlinebreaklocale "zh"
\renewcommand{\abstractname}{摘要}
\renewcommand{\contentsname}{目錄}
\renewcommand{\figurename}{圖}
\renewcommand{\tablename}{表}
\renewcommand{\refname}{參考文獻}
\begin{document}
\title{測量板的位置}\author{}\date{2026年5月14日}\maketitle
\begin{abstract}木板上有兩條刻線，用來記錄初始位置。\end{abstract}
\tableofcontents
\section{刻線}\begin{quote}「先讀下面的刻線。」\end{quote}
\begin{figure}[ht]\centering\fbox{\rule{0pt}{8mm}\rule{33mm}{0pt}}\caption{木板位置}\end{figure}
\section{高度}
\begin{table}[ht]\centering\begin{tabular}{cr}刻線&高度 (mm)\\\hline 1&12\\2&18\end{tabular}\caption{刻線高度}\end{table}
\begin{thebibliography}{1}\bibitem{record}\emph{測量手冊}, 2025, 第2章.\end{thebibliography}
\end{document}
